\section{Round-five positive closure: recession-complete collision path large deviations}
\label{sec:r5-a3}

The inducing scheme is completed by retaining, rather than discarding, the
empirical mass of long excursions.  The lower bound uses entropy-preserving
finite-state Markov approximants; periodic orbit measures are not used to
approximate positive-entropy laws.

\subsection{The countable graph completion}

Paper A2 supplies a fixed countable cut tower for the collision map.  Denote
its one-sided symbolic graph by \(\mathcal G_R=(\mathcal A_R,E_R)\), its shift
by \(\sigma\), and the coding map by \(\pi_R\).  The alphabet records the
homogeneity strip, lifted disk, one-sided singular germ and return level.  Put
\[
 d_\vartheta(x,y)=\vartheta^{n(x,y)},
 \qquad W(x)=1+r(x)+|k(x)|,
\]
where \(r\) is the return level and \(k\) the lifted displacement.  The A2
complexity estimates imply
\[
 \sum_{a\in\mathcal A_R}e^{-cW(a)}<\infty
\]
with one \(c>0\), uniformly in \(R\).

Choose a finite union \(Y_R\) of magnet rectangles.  Let
\(\mathcal G_R^{\rm ret}\) be the graph of paths which return to \(Y_R\), and
let \(\mathcal G_R^{\infty}\) be the closed survivor subgraph obtained by
deleting all vertices in \(Y_R\).  No finite-depth truncation is declared
forward invariant.  All survivor paths remain in the complete countable
subgraph.

For a bounded dynamically H\"older path potential \(F\), let
\(\Phi_{R,F}\) be its one-step lift.  On the weighted variation space
\[
 \mathscr V_\beta=\left\{h:
 \|h\|_\beta=\sup_a e^{-\beta W(a)}|h|_{[a]}
 +\sum_{n\ge1}\vartheta^{-n}\operatorname{var}_n(h)<\infty
 \right\}
\]
define the recurrent and survivor Ruelle operators
\[
 \mathcal R_{R,F}(z)h
 =\sum_{\substack{\sigma y=x\\y\text{ first return}}}
 e^{S_{r(y)}\Phi_{R,F}(y)-zr(y)}h(y),
\]
\[
 \mathcal S_{R,F}h(x)
 =\sum_{\substack{\sigma y=x\\y\in\mathcal G_R^\infty}}
 e^{\Phi_{R,F}(y)}h(y).
\]
The first is a renewal operator on the magnet and the second is the genuine
open-system operator on the full survivor graph.

\begin{lemma}[Common operator domains and renewal decomposition]
\label{lem:r5-a3-renewal}
There are \(\beta>0\) and \(\gamma>0\), uniform on compact source sets, such
that:
\begin{enumerate}[label=(\roman*)]
\item \(\mathcal R_{R,F}(z)\) is analytic on
      \(\Re z>q_Y(F)-\gamma\) as a bounded operator on
      \(\mathscr V_\beta(Y_R)\);
\item \(\mathcal S_{R,F}\) is quasi-compact on the survivor weighted space;
\item the entrance and exit operators between the two spaces are bounded in
      the same half-strip; and
\item the full collision resolvent has the exact block factorization
\[
 \mathscr G_{R,F}(z)=
 \mathscr E^{\rm in}_{R,F}(z)
 (I-\mathcal R_{R,F}(z))^{-1}
 \mathscr E^{\rm out}_{R,F}(z)
 +\mathscr G^{\infty}_{R,F}(z).
\]
\end{enumerate}
Every term acts between the explicitly declared weighted spaces.
\end{lemma}

\begin{proof}
Decompose a coded path at its successive visits to the magnet.  The A2 tower
weight gives
\[
 \sum_{r=n}\sup_{[a]}e^{S_n\Re\Phi_{R,F}}\le Ce^{-\gamma n}
\]
after reducing \(\gamma\) on a compact source set.  This proves absolute
operator convergence of the return, entrance and exit series in the stated
half-strip.  On the survivor graph the same weight gives a Doeblin--Fortet
inequality: variations contract by \(\vartheta\), while the weighted
supremum tail is compact.  Hence the survivor operator is quasi-compact.
Every path is uniquely either an infinite survivor or a finite entrance,
zero or more complete return blocks, and an exit.  Multiplying the
corresponding operator weights and summing the geometric series proves the
factorization.  Graph-completed singular paths are assigned by their oriented
one-sided code, so no boundary orbit is omitted or counted twice.
\end{proof}

\subsection{Direct pressure with component labels}

Let \(q_Y(F)\) be the unique real number at which the leading eigenvalue of
\(\mathcal R_{R,F}(z)\) equals one.  The survivor graph has at most countably
many topologically transitive components \(D\).  Put
\[
 q_\infty(F)=\sup_D P_D(\Phi_{R,F}),
 \qquad Q_R^c(F)=\max\{q_Y(F),q_\infty(F)\}.
\]
A survivor component label is retained in the phase state; dominance over the
return branch is not confused with uniqueness inside the survivor sector.

\begin{theorem}[Recurrent--survivor pressure theorem]
\label{thm:r5-a3-pressure}
The collision-time logarithmic moment generating function exists and equals
\(Q_R^c(F)\).  On a chart where either the recurrent root or one survivor
component is separated by a positive pressure gap, the pressure and its
unique equilibrium phase are analytic.  At a tie the maximizing phase set is
the compact convex hull of the finitely many components within the prescribed
pressure tolerance; the pressure remains convex and continuous.
\end{theorem}

\begin{proof}
The abscissa of convergence of the first term in
Lemma~\ref{lem:r5-a3-renewal} is \(q_Y(F)\).  The second term has abscissa
\(q_\infty(F)\) by the spectral radius formula on the survivor components.
The logarithmic growth of their sum is the maximum.  A positive gap isolates
a simple eigenvalue on one declared component; analytic perturbation gives the
first assertion.  The weighted graph has only finitely many components with
pressure above \(Q_R^c(F)-\delta\), because each such component must meet the
finite weight sublevel \(\{W\le C_\delta\}\).  This proves compactness and the
finite convex-hull description at coexistence.
\end{proof}

\subsection{Entropy-preserving Markov approximation}

Let \(\mathcal M_R\) be the invariant probability measures on the full coded
collision path with finite \(W\)-moment.  For \(\nu\in\mathcal M_R\), write
\(\mathcal P_m\) for the finite partition generated by all cylinders whose
symbols satisfy \(W\le m\), together with one tail atom.  Let
\(p_m(B'\mid B)\) be the conditional transition probabilities of \(\nu\)
on the \(m\)-block states.  Replacing zero transitions by
\(\delta_m e^{-W(B')}\) and renormalizing gives an irreducible finite-state
Markov law \(\nu_m\).

\begin{lemma}[Rate-dense exposed Markov phases]
\label{lem:r5-a3-density}
For every invariant \(\nu\) with finite pressure-dual cost there are
finite-state irreducible Markov laws \(\nu_m\) such that
\[
 \nu_m\Rightarrow\nu,\qquad
 \nu_m(W)\to\nu(W),\qquad
 h(\nu_m)\to h(\nu),\qquad
 \nu_m(F)\to\nu(F)
\]
for every source in a fixed compact weighted set.  Each \(\nu_m\) is the
unique equilibrium state of the locally constant exposing potential
\[
 G_m(x)=\log p_m(x_1\mid x_0)-\log Z_m(x_0).
\]
The construction works inside the survivor subgraph when \(\nu(Y_R)=0\).
\end{lemma}

\begin{proof}
The tail atom has \(\nu\)-mass and \(W\)-moment tending to zero.  Martingale
convergence of the finite conditional probabilities gives weak convergence of
the Markovizations.  Entropy is the decreasing limit of conditional
entropies,
\[
 h(\nu)=\lim_{m\to\infty}
 H_\nu(\mathcal P_m\mid\mathcal P_m^-).
\]
The entropy of \(\nu_m\) is exactly the corresponding finite conditional
entropy up to an error
\(O(\delta_m|\log\delta_m|)\); choose \(\delta_m\downarrow0\) after the tail
cutoff.  This proves entropy convergence and, by weighted uniform
integrability, convergence of every declared potential.  The Perron--Frobenius
theorem for the finite irreducible transition matrix makes \(\nu_m\) the
unique equilibrium of \(G_m\).  If \(\nu(Y_R)=0\), choose only survivor
cylinders; the same construction never introduces a magnet transition.
Periodic measures play no role.
\end{proof}

\subsection{A diagonal singularity shield}

Let \(B_{q,\delta}\) be a smooth indicator of points whose first \(q\)
iterates approach a singularity closer than \(\delta\).  The growth lemma and
the one-step homogeneity estimate give constants \(c_0,c_1>0\) such that
\[
 Q_R^c(sB_{q,\delta})-Q_R^c(0)
 \le C e^{c_0s}\bigl(e^{-c_1q}+q^C\delta^\alpha\bigr).
\]
For large \(s\), choose
\[
 q(s)=\lceil 2c_0s/c_1\rceil,
 \qquad \delta(s)=e^{-4c_0s/\alpha}q(s)^{-2C/\alpha}.
\]
Then the right-hand side tends to zero while the reward on the singular set is
\(s\).

\begin{lemma}[Infinite cost of unresolved singular mass]
\label{lem:r5-a3-shield}
Every path law charging the unresolved singular set has infinite pressure-dual
cost.  Regular cylinder potentials are therefore exponentially good
approximations of graph-completed path observables.
\end{lemma}

\begin{proof}
The pressure estimate follows by decomposing singular visits into maximal
clusters.  A cluster of length \(q\) has standard-family weight at most
\(Ce^{-c_1q}+Cq^C\delta^\alpha\); selecting clusters contributes at most the
exponential generating factor \(e^{c_0s}\).  With the diagonal choice above,
Fenchel duality gives for a measure with singular mass \(m>0\)
\[
 I(\nu)\ge sm-o(1).
\]
Letting \(s\to\infty\) proves the claim.  Notice that no pressure theorem at a
fixed singular neighborhood and arbitrarily large source is assumed.
\end{proof}

\subsection{The full collision empirical-path LDP}

Let
\[
 L_n^c=\frac1n\sum_{j=0}^{n-1}\delta_{\Theta^j\omega}
\]
on the graph-completed collision path space.  Put
\[
 I_R^c(\nu)=
 \begin{cases}
  Q_R^c(0)-h_\nu(\Theta)-\nu(\Phi_{R,0}),
  &\nu\in\mathcal M_R,\\
  +\infty,&\text{otherwise}.
 \end{cases}
\]
Equivalently,
\[
 I_R^c(\nu)=\sup_F\{\nu(F)-Q_R^c(F)+Q_R^c(0)\}
\]
on the weighted strict source domain.

\begin{theorem}[Full collision empirical-path LDP]
\label{thm:r5-a3-collision-ldp}
The laws of \(L_n^c\) satisfy a full good LDP with speed \(n\) and rate
\(I_R^c\), uniformly for \(R\) in compact subintervals of \(\mathcal I\).
The lower bound holds at every finite-rate invariant law, including
positive-entropy survivor laws and laws with a nonzero recession fraction of
long excursions.
\end{theorem}

\begin{proof}
For a finite alphabet/weight truncation the graph is a finite union of Markov
components.  The usual Perron change of law gives the level-three upper and
lower bounds.  Weighted exponential tightness follows from the tower tail and
Lemma~\ref{lem:r5-a3-shield}.  Hence the projective upper bound passes to the
countable graph and has compact sublevels.

For the lower bound, first take a finite-state irreducible Markov phase.  Its
exposing potential gives an exact change of measure; Shannon--McMillan and the
Markov law of large numbers yield the required open-set exponent.  For a
general finite-rate \(\nu\), use the laws \(\nu_m\) from
Lemma~\ref{lem:r5-a3-density}.  Their entropies, energies and weighted moments
converge, so their lower-bound costs converge to \(I_R^c(\nu)\).  Exponential
tightness permits the diagonal limit.  Long excursions are represented by
states with increasing return level; their limiting mass is retained by the
weighted topology and is not discarded as a terminal remainder.
\end{proof}

\subsection{Physical time and terminal clock surgery}

Let \(\tau_R\) be the bounded finite-horizon roof and define the Palm map
\[
 \mathfrak P_R(\nu)(A)=
 \frac1{\nu(\tau_R)}
 \int\int_0^{\tau_R(\omega)}
 1_A(\Theta_t\omega)\,dt\,d\nu(\omega).
\]
On the graph completion the one-sided roof values are separate coordinates,
so this map is continuous on every finite-rate sublevel.  The last incomplete
flight has length at most \(\tau_+\); replacing it by a complete rooted flight
changes a bounded path average by \(O(T^{-1})\).  This is terminal clock
surgery, not an exponential truncation of a macroscopic return block.

\begin{theorem}[Full physical-time empirical-path LDP]
\label{thm:r5-a3-physical-ldp}
The physical empirical path measures satisfy a good LDP with rate
\[
 I_R^{\rm ph}(\rho)=
 \inf\left\{
  \frac{I_R^c(\nu)}{\nu(\tau_R)}:
  \mathfrak P_R(\nu)=\rho
 \right\}.
\]
The collision and physical formulations are linked by the continuous Palm map
and the bounded terminal surgery above.
\end{theorem}

\begin{proof}
Finite horizon gives
\(0<\tau_-\le\nu(\tau_R)\le\tau_+\) on invariant laws.  Continuity of the
Palm map and the contraction principle give the upper bound and goodness.
For the lower bound, use the Markov approximants from
Lemma~\ref{lem:r5-a3-density}; their roof means and Palm images converge.
Concatenate complete blocks and perform the bounded terminal surgery.  Its
cost and path error vanish after division by physical time, yielding the
displayed contraction rate.
\end{proof}

\subsection{Conditioning and Hausdorff cotangents}

Let \(C\) be a continuous finite-dimensional observable and let
\(a\) be in the relative interior of its finite-rate image.  For balls
\(B(a,\delta_n)\), choose \(\delta_n\downarrow0\) slowly enough that the LDP
lower bound for one exposed approximant is uniform; equivalently
\(n^{-1}\log\mathbb P\{C(L_n)\in B(a,\delta_n)\}\) converges to the constrained
minimum.  This condition is automatic for fixed neighborhoods and can always
be met by a diagonal sequence.

Let \(C_W\) be the weighted-strict potential space and let
\[
 \mathscr N_W=\overline{\mathbb R\mathbf1+
 \{U-U\circ\Theta_t:t\in\mathbb Q,\ U\in C_W\}}^{\,\beta_W}.
\]

\begin{theorem}[Information projection and cotangent quotient]
\label{thm:r5-a3-cotangent}
The conditional empirical laws concentrate exponentially on the minimizer set
of \(I_R^c\) under \(C(\nu)=a\).  If the minimizer is unique the rooted path
law converges to it; otherwise every subsequential conditional law is a
probability mixture supported on the compact minimizer set.
Furthermore
\[
 C_W/\mathscr N_W
\]
is Hausdorff, and two potentials define the same tangent on every invariant
finite-weight phase exactly when their difference belongs to
\(\mathscr N_W\).
\end{theorem}

\begin{proof}
Goodness gives a compact minimizer set.  The LDP upper bound outside a
neighborhood of this set and the matching denominator lower exponent give
exponential conditional concentration.  Tightness of rooted conditional laws
follows from the same rate sublevel; disintegration then gives the mixture
statement.

The space \(\mathscr N_W\) is closed by definition.  Its annihilator consists
of signed invariant Radon measures of total mass zero.  If a potential is not
in \(\mathscr N_W\), Hahn--Banach supplies such a separator \(m\).  Write the
Jordan decomposition and add a fixed invariant finite-weight probability
\(\rho\) so that
\[
 \nu_\pm=\frac{m_\pm+c\rho}{m_\pm(1)+c}
\]
have equal normalizing mass.  Invariance of \(m\) makes \(m_+(1)=m_-(1)\),
so \(\nu_+\) and \(\nu_-\) are invariant probabilities whose integrals of the
potential differ.  This proves the converse annihilator assertion.
\end{proof}

\begin{theorem}[Round-five A3 closure]
\label{thm:r5-a3-main}
The collision and physical empirical path measures satisfy full good LDPs
which retain survivor and long-excursion mass.  The lower bound is supplied by
entropy-preserving exposed Markov phases, singular paths have infinite cost,
regular information projections obey Gibbs conditioning, and the cotangent is
the declared weighted-strict Hausdorff quotient.
\end{theorem}

\begin{proof}
Combine Theorems~\ref{thm:r5-a3-pressure},
\ref{thm:r5-a3-collision-ldp}, \ref{thm:r5-a3-physical-ldp}, and
\ref{thm:r5-a3-cotangent}.
\end{proof}
